% =============================================================================
% Section VII: Discussion
% paper/discussion.tex
% =============================================================================

\section{Discussion}
\label{sec:discussion}

The results of \cref{sec:powerlaw,sec:desitter,sec:em} establish that causal
horizons in FLRW spacetimes produce inevitable, unbounded decoherence of quantum
superpositions of charged bodies, even in the complete absence of Killing
symmetry.  In this section we examine the physical mechanism underlying this
result, explain the geometric origin of the two qualitatively different growth
laws, discuss the universality of the factor-of-2 relation between scalar and
electromagnetic decoherence, compare our findings to related work in the
literature, collect all results in a single summary table, and identify
directions for future investigation.

\subsection{Physical mechanism: soft radiation through the horizon}
\label{sec:discussion-mechanism}

The decoherence computed in this paper has a transparent physical origin.
A body in a quantum superposition of two spatial locations $\gamma_1$ and
$\gamma_2$ sources a field whose radiation pattern encodes which-path information:
the field emitted by a charge at $\gamma_1$ differs from that emitted by a charge
at $\gamma_2$ precisely in the dipole mode.  This difference constitutes soft
radiation that propagates outward from the superposition.  When a future causal
horizon $\mc{H}^+(\gamma)$ is present, a portion of this radiation — specifically,
the low-frequency (long-wavelength) modes emitted during the superposition — crosses
the horizon and becomes permanently inaccessible to any comoving observer.  In the
DSW language, the in-vacuum field is in an entangled state across the horizon; after
the radiation has crossed, the reduced state of the observable field is mixed, and
tracing over the inaccessible modes suppresses the off-diagonal coherence
$|\langle \gamma_1 | \rho_{\rm red} | \gamma_2 \rangle|$.

The technical manifestation of this mechanism is the role of the conformal time
range as an infrared (IR) regulator.  The decoherence integral,
\cref{eq:decoherence-integral}, involves a double integral over conformal
time with an integrand $\sim (t_1 - t_2)^{-4}$.  In a Minkowski background there
is no upper bound on the conformal time, and the integral, evaluated over an
unbounded interval, yields a contribution that is entirely insensitive to the total
integration time at long wavelengths.  By contrast, when a future causal horizon
is present ($p > 1$ or de Sitter), the conformal time range is bounded above by
the finite value $t_\infty < \infty$.  The horizon imposes a physical IR cutoff on
the field modes that contribute to decoherence: modes with wavelengths longer than
the remaining conformal distance $t_\infty - t$ cannot fit between the source and
the horizon and are therefore not available to carry which-path information.  As
the superposition time $T$ increases and $t_{\rm off} \to t_\infty$, progressively
shorter-wavelength (higher-frequency) modes become saturated, and the cumulative
decoherence grows without bound.

The $p = 1$ case makes this causal interpretation precise.  When $p = 1$ the
conformal time range is all of $[t_0, \infty)$: no future causal horizon is
present, and there is no IR cutoff.  The integral in \cref{eq:decoherence-integral}
then yields a finite constant as $T \to \infty$ (see \cref{eq:N-nohorizon}): the
soft radiation is \emph{displaced} into the field at spatial infinity but is not
\emph{permanently hidden} from any observer.  The horizon, not the background
dynamics, is the essential ingredient.

\subsection{Logarithmic vs.\ linear growth: a geometric distinction}
\label{sec:discussion-growth}

The two growth laws found in this paper — $\Nnum \sim \ln T$ for power-law FLRW
and $\Nnum \sim T$ for de Sitter — reflect a geometric difference in how the
respective horizons ``peel off'' soft field modes.

In de Sitter spacetime, the exponential scale factor $\Scfac(\eta) = -1/(H\eta)$
means that the conformal-time distance from the source to the future boundary
shrinks in inverse proportion to $-\eta$ as proper time advances.  New modes
cross the horizon at a \emph{constant} proper-time rate.  The Wightman function
at coincident spatial position, \cref{eq:Wight-dS-coincident}, acquires the
uniform spectral weight $H^2$, which is time-independent: every moment in
de Sitter looks identical to every other, because the spacetime is maximally
symmetric.  The decoherence integral then accumulates linearly in proper time,
and the rate is set by the Gibbons-Hawking temperature $T_{\rm GH} = H/2\pi$
through the identification $H^3 = (2\pi T_{\rm GH})^3$.  From the perspective
of a comoving observer, decoherence arises from coupling to a genuine thermal
bath at temperature $T_{\rm GH}$~\cite{Gibbons:1977mu,Danielson:2022tdw}; the
rate is exactly what a dipole of moment $qd_0$ accumulates when immersed in a
Planck distribution of soft bosons~\cite{Danielson:2024yge}.

For the power-law family with $p > 1$, the expansion is accelerated but
\emph{not} exponential: the scale factor $\Scfac(\tau) = C_p(\tau - \tau_*)^p$
grows as a power law.  The conformal-time distance remaining to the horizon
decreases, but at a rate that slows as proper time advances.  Each successive
interval $d\tau$ of superposition covers a smaller conformal-time fraction of
the remaining horizon distance.  The cumulative decoherence, which counts the
number of modes that have crossed the horizon, therefore grows as the logarithm
of the total proper time elapsed.  There is no Killing symmetry, no invariant
notion of temperature, and no thermal distribution: the decoherence rate is not
set by a bath temperature but purely by the geometry of the horizon, encoded in
the single parameter $T_0$ defined by \cref{eq:T0-def}.

Both laws are unbounded — decoherence is inevitable in both cases — but they
are parametrically distinct.  A superposition that survives for proper time $T$
accumulates $\Nnum_{\rm dS} \sim H^3 T$ (linear) or
$\Nnum_{\rm PL} \sim p^3 T_0^{-2} \ln(T/T_0)$ (logarithmic).  For any fixed
$q$, $d_0$, and background parameters, the de Sitter rate dominates at late
times.  Conversely, the $p^3$ prefactor in the power-law result shows that
stronger acceleration (larger $p$, shorter horizon distance) produces faster
logarithmic growth: an observer in a more violently accelerating cosmology
loses coherence more rapidly even within the slower logarithmic regime.

\subsection{Universality: conformal invariance and the factor of 2}
\label{sec:discussion-universality}

The equality $\Nnum_{\rm EM} = 2\,\Nnum_{\rm scalar}$, derived in \cref{sec:em},
holds for every member of the power-law family and for de Sitter, identically.
This universality has a clean structural explanation.

Both the conformally coupled scalar and the Maxwell field are conformally
invariant in four spacetime dimensions.  Under the Weyl rescaling
$g_{ab} \to \Scfac^2(\eta)\,\eta_{ab}$, the equations of motion for both fields
reduce to their flat-space counterparts: the scalar field equation with its
$R/6$ coupling becomes the massless Klein-Gordon equation, and the Maxwell
equations in Lorenz gauge become the flat-space wave equation (see
\cref{sec:em-conformal}).  Consequently, the structure of the decoherence
integral is identical for both fields: it is the \emph{flat-space} integral,
\cref{eq:decoherence-integral}, dressed only by the polarization content of the
field.  The scalar contributes one degree of freedom per propagation direction;
the photon contributes two independent transverse polarizations.  The factor of 2
counts these polarizations and is kinematic — independent of the cosmological
background, the power-law exponent, or any other geometric data.

This conformal universality extends to any conformally invariant field in any
conformally flat FLRW cosmology: the decoherence integral reduces to its Minkowski
counterpart, and the ratio of decoherence numbers between any two such fields is
fixed by their polarization content alone.  The results of this paper therefore
carry an intrinsic robustness: they do not depend on the detailed form of the
scale factor beyond the existence of a future causal horizon.

The gravitational case is qualitatively different.  Soft gravitons in four
spacetime dimensions are \emph{not} conformally invariant: the linearized
Einstein equations for the metric perturbation $h_{\mu\nu}$ do not reduce to
the flat-space wave equation under a Weyl rescaling.  The conformal-flatness
argument used throughout this paper does not apply, and the calculation of
gravitational decoherence in FLRW requires entirely new methods.  This is an
important open problem: in de Sitter, the gravitational decoherence number for
a massive body of mass $m$ and separation $d_0$ grows as
$\Nnum_{\rm grav}^{\rm dS} \sim H^5 m^2 d_0^4 T$~\cite{Danielson:2022tdw},
but the analogue for non-stationary FLRW cosmologies is not currently known.

\subsection{Comparison with related work}
\label{sec:discussion-comparison}

The de Sitter result~\cref{eq:N-dS},
$\Nnum_{\rm dS} = q^2 d_0^2 H^3 T/(96\pi^2)$, agrees exactly with the
result of Danielson, Satishchandran, and Wald~\cite{Danielson:2022tdw}, who
obtained the same expression by applying a Killing-horizon analysis in the static
patch of de Sitter.  Our derivation reaches this result by a completely different
route — conformal flatness and the FLRW Wightman function, with no reference to
the static patch or its Killing symmetry — and thereby confirms both calculations
and makes their physical equivalence explicit.  DSW (2025a)~\cite{Danielson:2024yge}
further showed that this result admits a local interpretation in terms of
low-frequency Hawking quanta at the observer's location; our calculation is
consistent with this local picture, which also requires Killing symmetry.

Li et al.~\cite{Li:2024} independently computed the decoherence of spatial
superpositions in de Sitter using algebraic quantum field theory methods.  Their
approach is complementary to ours: it works directly in the de Sitter algebra
rather than exploiting conformal flatness.  Their result for the scalar and
electromagnetic cases agrees with \cref{eq:N-dS,eq:N-EM-dS}, providing a
third independent confirmation of the numerical coefficient.

Gralla and Wei~\cite{Gralla:2023oya} developed a general formulation of
horizon-induced decoherence that covers Kerr black holes and other Killing-horizon
spacetimes.  Their framework is consistent with the DSW result for de Sitter; it
applies in the regime where a Killing horizon and a well-defined Hawking
temperature are available, which is precisely the regime where our power-law
formula reduces to linear growth.  The present paper extends beyond the Killing
symmetry assumption: our main new result — the logarithmic growth law
\cref{eq:N-large-T} — has no analogue in the Gralla-Wei framework, which requires
stationarity.

Wilson-Gerow, Dugad, and Chen~\cite{WilsonGerow:2024} analyzed horizon-induced
decoherence through the Unruh-DeWitt lens, treating the horizon as a ``warm
bath'' at the Gibbons-Hawking temperature.  Their results are consistent with the
thermal ($\Nnum \sim T$) growth for Killing horizons.  For non-stationary FLRW
backgrounds, no Unruh temperature is defined, and the warm-bath approximation
does not apply; the logarithmic growth found in this paper lies outside their
framework.

Fahn and Pesci~\cite{FahnPesci:2025} recently studied quantum geometry
corrections to the DSW decoherence near a black hole horizon, finding that
Planck-scale discretization of the horizon area can \emph{suppress} decoherence
relative to the semiclassical prediction at very short distances.  The results
of the present paper are the semiclassical baseline against which such
corrections would be measured.  Extending the quantum-geometry analysis to
non-stationary FLRW horizons is an interesting open problem.

\subsection{Summary of results}
\label{sec:discussion-table}

\Cref{tab:summary} collects the decoherence number $\Nnum$ for all cases studied
in this paper.  The electromagnetic result $\Nnum_{\rm EM} = 2\,\Nnum_{\rm scalar}$
holds in every row.

\begin{table}[h]
\caption{Decoherence number $\Nnum$ for a scalar field superposition of charge $q$,
separation $d_0$, and duration $T$.  The parameter $T_0$ is the characteristic
proper-time scale defined in \cref{eq:T0-def}, and $\widetilde{T} \equiv T/T_0$.
The electromagnetic result is $\Nnum_{\rm EM} = 2\,\Nnum_{\rm scalar}$
in all cases.  The exact closed-form expression for power-law FLRW
is \cref{eq:N-exact}; the large-$T$ asymptotic is quoted here.}
\label{tab:summary}
\begin{ruledtabular}
\begin{tabular}{lccc}
Spacetime & Scale factor & $\Nnum\;(T \to \infty)$ & Growth \\[2pt]
\hline\\[-6pt]
Power-law, $p > 1$ &
  $C_p(\tau - \tau_*)^p$ &
  $\dfrac{q^2 d_0^2 p^3}{48\pi^3 T_0^2}\ln\!\lb(\dfrac{T}{T_0}\rb)$ &
  logarithmic \\[14pt]
No horizon, $p = 1$ &
  $C_1(\tau - \tau_*)$ &
  $\dfrac{q^2 d_0^2}{16\pi^2 u_{\rm on}^2}$ &
  saturates \\[14pt]
De~Sitter &
  $e^{H\tau}/H$ &
  $\dfrac{q^2 d_0^2 H^3}{96\pi^2}\, T$ &
  linear \\[4pt]
\end{tabular}
\end{ruledtabular}
\end{table}

The table illustrates the central hierarchy: the causal horizon converts a
saturating finite constant into an unbounded decoherence number.  Among
horizon-equipped spacetimes, the growth rate depends on whether the horizon
is stationary (linear, de Sitter) or non-stationary (logarithmic, power-law).

\subsection{Implications and future directions}
\label{sec:discussion-future}

\paragraph{Observational scales.}
For a superposition of a charged body with charge $q$, comoving separation $d_0$,
and initial conformal-time distance $T_0$ to the causal horizon, the decoherence
number \cref{eq:N-large-T} implies a characteristic coherence time: requiring
$\Nnum \sim 1$ gives
%
\be
  T_{\rm dec}
  \;\sim\;
  T_0\, \exp\!\lb(\frac{48\pi^3\, T_0^2}{q^2 d_0^2 p^3}\rb) \,.
  \label{eq:T-dec}
\ee
%
For any macroscopic charged body ($q$ of order the elementary charge, $d_0$ of
order nanometers) in a cosmological background ($T_0 \sim H_0^{-1}$ for the
present Hubble horizon), $T_{\rm dec}$ is astronomically large — far exceeding
the current age of the universe.  Nonetheless, \cref{eq:T-dec} is finite: given
sufficient time, any such superposition is completely decohered by the causal
horizon.  The result is therefore a statement of principle as much as practice:
\emph{no superposition is protected by the absence of a thermal bath, because the
horizon itself acts as an irreversible information sink}.  For small charges or
separations, or in spacetimes with large acceleration ($p \gg 1$), the
decoherence can be parametrically faster.

\paragraph{Gravitational decoherence in FLRW.}
The most important open problem is the extension to soft gravitons.  Unlike
conformally invariant matter fields, the linearized gravitational field is not
conformally invariant in four dimensions: the minimal coupling to curvature
cannot be removed by a Weyl rescaling, and the graviton Wightman function in
FLRW does not reduce to the Minkowski form.  A calculation of gravitational
decoherence for power-law FLRW cosmologies therefore requires a direct
evaluation of the graviton 2-point function in these backgrounds, likely
involving a more involved mode analysis.  This is a natural next step, as it
would determine the decoherence of massive (but uncharged) bodies in
cosmological spacetimes.

\paragraph{Realistic cosmology.}
The universe underwent a transition from matter domination ($p = 2/3$, no future
horizon) to $\Lambda$-domination (approximately de Sitter, with a causal horizon)
at a redshift $z_\Lambda \approx 0.3$.  A complete calculation of horizon
decoherence in a realistic FLRW model would need to track this transition and
identify when the IR cutoff provided by the horizon first becomes operative.  The
framework developed here — in particular the conformal-time formulation — provides
the natural starting point for such a computation.

\paragraph{Quantum geometry corrections.}
Fahn and Pesci~\cite{FahnPesci:2025} have shown that quantizing the horizon
area in the loop quantum gravity framework suppresses the semiclassical
decoherence by a factor that depends on the ratio of the separation $d_0$ to the
Planck length.  For FLRW horizons, the quantum geometry of the cosmological
horizon is less well understood than that of a black hole horizon; determining
whether analogous suppression occurs, and at what scale it becomes relevant, is
an important intersection of this work with approaches to quantum gravity.

\paragraph{Observable signatures.}
Both growth laws — logarithmic and linear in $T$ — give distinctive
\emph{time-dependence} to the coherence amplitude $|\langle \gamma_1 | \rho_{\rm
red} | \gamma_2\rangle| = e^{-\Nnum(T)}$.  Whether experiments sensitive to
the time profile of decoherence could distinguish horizon-induced decay from
other environmental decoherence mechanisms (thermal phonons, electromagnetic
fluctuations, etc.) is an open and practically motivated question, particularly
in light of proposals~\cite{Carney:2019} for tabletop tests of quantum gravity
phenomenology.

%% CITATIONS NEEDED (for gpd-bibliographer):
%% - Danielson:2022tdw  -- DSW (2023) PRD 108, 025007 (de Sitter N~H^3 T and gravitational result)
%% - Danielson:2024yge  -- DSW (2025a) PRD 111, 025014 (local Hawking-quanta interpretation)
%% - Gibbons:1977mu     -- Gibbons & Hawking (1977) PRD 15, 2738 (T_GH = H/2pi)
%% - Li:2024            -- Li et al. (2025) arXiv:2501.00213 (algebraic QFT de Sitter comparison)
%% - Gralla:2023oya     -- Gralla & Wei (2024) PRD 109, 065031 (Kerr black holes, general formulation)
%% - WilsonGerow:2024   -- Wilson-Gerow, Dugad & Chen (2024) PRD 110, 045002 (warm horizons)
%% - FahnPesci:2025     -- Fahn & Pesci (2025) PRD, arXiv:2507.16911 (quantum geometry suppression)
%% - Carney:2019        -- Carney, Stamp & Taylor (2019) CQG 36, 034001 (tabletop quantum gravity)
